\section{Preliminaries}

\subsection{Pauli--Liouville Representation}

To facilitate the analysis of distillation, we present the Pauli--Liouville representation \cite{Lin2021} in the two-qubit scenario below.
Let
\(
\mathbb{P}_2 \coloneqq\{I,X,Y,Z\}^{ \otimes 2}
\)
denote the two-qubit Pauli operator basis, which is a
complete orthogonal basis for the operator space
$\mathcal{L}\bigl((\mathbb{C}^{2})^{\otimes 2}\bigr)$ under the
Hilbert--Schmidt inner product. Then, every two-qubit operator $\mathcal{G}$ can be uniquely expressed as a linear combination of the elements of $\mathbb P_2$.
% admits a unique Pauli
% expansion as a linear combination of the Pauli basis elements.
Define the Pauli coefficient associated with
$P\in\mathbb{P}_2$ as
\begin{equation}
        s_P(\mathcal G)\coloneqq\operatorname{Tr}\!\left(P^{\dagger}\mathcal{G}\right),
\end{equation}
we can express $\mathcal{G}$ as follows,
\begin{equation}
    \mathcal{G}=\frac{1}{4}\sum_{P\in\mathbb{P}_2}s_P(\mathcal{G})P.
    \label{eq:pauli-expansion}
\end{equation}

Following the Pauli-Liouville representation,
we introduce the \(\left|\cdot\right\rangle\!\rangle\)  notation, and with it, the state can be expressed as a column vector as follows:
\begin{equation}
    \begin{aligned}
        \left|\rho\right\rangle\!\rangle
        \coloneqq\bigl[s_P(\rho)\bigr]^T_{P\in\mathbb{P}_2}.
    \end{aligned}
    \label{eq:pauli-liouville-vectors}
\end{equation}

For a two-qubit linear map \(\mathcal E\), its Pauli transfer matrix (PTM) \(R_{\mathcal E}\) is a \(16\times16\) matrix, whose enties are
\begin{equation}
    \left(R_{\mathcal E}\right)_{P,Q}
    \coloneqq\frac{1}{4}\Tr\!\left[P^\dagger\mathcal E(Q)\right],
    \text{ for any } P,Q\in\mathbb{P}_2.
    \label{eq:ptm-definition}
\end{equation}
Consequently, the outcome state of applying the operator \(\mathcal E\) to an input state \(\rho\) can be expressed as follows,
\begin{equation}
    \begin{aligned}
        \left|\mathcal E(\rho)\right\rangle\!\rangle
         & =R_{\mathcal E}\left|\rho\right\rangle\!\rangle.
    \end{aligned}
    \label{eq:pauli-liouville-born-rule}
\end{equation}


\subsection{Pauli Channels in Bell-Diagonal Basis}

% Bell states are particularly convenient for describing two-qubit
% entanglement under Pauli noise because single-qubit Pauli errors
% permute the Bell states. 
Bell basis (composed by Bell states) is particularly convenient for describing two-qubit
entanglement under Pauli noise because applying a single-qubit Pauli error to one half of a Bell pair maps it to another Bell state.
The four Bell states are
\begin{equation}
    \ket{\Phi^\pm}
    =
    \frac{\ket{00}\pm\ket{11}}{\sqrt{2}},
    \qquad
    \ket{\Psi^\pm}
    =
    \frac{\ket{01}\pm\ket{10}}{\sqrt{2}}.
\end{equation}

Their density operators admit the Pauli decomposition
\begin{equation}
    \begin{aligned}
        \ket{\Phi^+}\!\bra{\Phi^+}
         & =
        \frac{1}{4}
        \left(
        I I
        +X X
        -Y Y
        +Z Z
        \right), \\
        \ket{\Psi^+}\!\bra{\Psi^+}
         & =
        \frac{1}{4}
        \left(
        I I
        +X X
        +Y Y
        -Z Z
        \right), \\
        \ket{\Psi^-}\!\bra{\Psi^-}
         & =
        \frac{1}{4}
        \left(
        I I
        -X X
        -Y Y
        -Z Z
        \right), \\
        \ket{\Phi^-}\!\bra{\Phi^-}
         & =
        \frac{1}{4}
        \left(
        I I
        -X X
        +Y Y
        +Z Z
        \right).
    \end{aligned}
    \label{eq:bell-projector-pauli-expansions}
\end{equation}

These decompositions motivate the definition of the operator subspace
\begin{equation}
    \mathcal S_{\mathrm B}
    \coloneqq
    \operatorname{span}
    \{I I, X X, Y Y, Z Z\}.
    \label{eq:bell-diagonal-subspace}
\end{equation}
An operator belongs to $\mathcal S_{\mathrm B}$ if it is a linear
combination of these four Pauli operators.

A two-qubit density operator $\rho$ is called Bell-diagonal if it
is diagonal in the Bell basis. By
Lemma~\ref{lem:bell-pauli-transformation}, this is equivalent to
$\rho\in\mathcal S_{\mathrm B}$.

For such a state, we define its restricted Pauli--Liouville vector as
\begin{equation}
    \plket{\rho}
    \coloneqq
    \begin{bmatrix}
        s_{II}(\rho)  &
        s_{XX}(\rho)  &
        s_{Y Y}(\rho) &
        s_{Z Z}(\rho)
    \end{bmatrix}^{ T},
    \label{eq:restricted-pauli-liouville-vector}
\end{equation}
where $s_P(\rho)\coloneqq\operatorname{Tr}(P^\dagger\rho)$.

The relation between the Bell basis and the Pauli basis is
standard~\cite{Horodecki1996,Luo2008}. For completeness of presentation, we state the
corresponding change-of-basis relation below.

\begin{lemma}
    \label{lem:bell-pauli-transformation}
    Consider the following ordered bases of \(\mathcal S_B\):
    \begin{equation}
        \begin{aligned}
            \mathcal B_{\mathrm{Bell}}
             & \coloneqq
            \left(
            \ket{\Phi^+}\!\bra{\Phi^+},
            \ket{\Psi^+}\!\bra{\Psi^+},
            \ket{\Psi^-}\!\bra{\Psi^-},
            \ket{\Phi^-}\!\bra{\Phi^-}
            \right),     \\
            \mathcal B_{\mathrm P}
             & \coloneqq
            \left(
            I I,
            X X,
            Y Y,
            Z Z
            \right).
        \end{aligned}
        \label{1}
    \end{equation}
    There exists a full-rank transformation matrix $W$ which satisfies $4WW^T= I_4$ such that
    \begin{equation}
        \mathcal{B}_{\mathrm{Bell}}
        =
        \mathcal{B}_{\mathrm{P}}W,
        \qquad
        W=
        \frac{1}{4}     \begin{bmatrix*}[r]
         1 &  1 &  1 &  1 \\
         1 &  1 & -1 & -1 \\
        -1 &  1 & -1 &  1 \\
         1 & -1 & -1 &  1
    \end{bmatrix*}.
        \label{eq:bell-pauli-transform}
    \end{equation}
\end{lemma}

For a normalized state \(\rho\), its fidelity is defined as
\begin{equation}
    \begin{aligned}
        F(\rho)
         & =\frac{1}{4}\left[
                           s_{I I}(\rho)
                           +s_{X X}(\rho)
                           -s_{Y Y}(\rho)
                           +s_{Z Z}(\rho)
                           \right].
    \end{aligned}
    \label{eq:bell-state-fidelity}
\end{equation}

Consider the single-qubit Pauli channel
\begin{equation}
    \mathcal P(\sigma)
    =
    A\sigma
    +BX\sigma X
    +CY\sigma Y
    +DZ\sigma Z,
    \label{eq:single-qubit-pauli-channel}
\end{equation}
where
\(
A, B, C, D\geq 0
\) and \(
A+B+C+D=1.
\)
If the first qubit is unaffected and the second qubit undergoes
\(\mathcal P\), the corresponding two-qubit channel is
\(
\mathcal E=\mathcal I\otimes\mathcal P,
\)
where \(\mathcal I\) denotes the identity channel.
Applying this channel to the maximally entangled state
\(\ket{\Phi^+}\) produces the Bell-diagonal state
\begin{equation}
    \begin{aligned}
        \rho
         & =
        \mathcal E\left(\ket{\Phi^+}\!\bra{\Phi^+}\right) \\
         & =
        A\ket{\Phi^+}\!\bra{\Phi^+}
                       {+}B\ket{\Psi^+}\!\bra{\Psi^+}
                                         {+}C\ket{\Psi^-}\!\bra{\Psi^-}
                                                           {+}D\ket{\Phi^-}\!\bra{\Phi^-}.
    \end{aligned}
    \label{eq:pauli-channel-output-state}
\end{equation}


As defined in \eqref{eq:ptm-definition}, the full PTM
\(R_{\mathcal E}\) is a \(16\times16\) matrix with respect to the
complete two-qubit Pauli basis. However, $\mathcal S_{\mathrm B}$ is invariant under $\mathcal E$,
meaning that $\mathcal E(\mathcal{G})\in\mathcal S_{\mathrm B}$ for every
$\mathcal{G}\in\mathcal S_{\mathrm B}$. Therefore, the action of \(\mathcal E\) on
Bell-diagonal states is completely described by its restriction to
\(\mathcal S_B\). With respect to the ordered basis
\(\mathcal B_{\mathrm P}\) in \ref{1}, the resulting \(4\times4\) PTM is
\begin{equation}
    \begin{aligned}
        \left.R_{\mathcal E}\right|_{\mathcal S_B}
         & =
        \operatorname{diag}(1,q_X,q_Y,q_Z), \\
        q_X
         & =
        A+B-C-D
        =
        2(A+B)-1,                           \\
        q_Y
         & =
        A-B+C-D
        =
        2(A+C)-1,                           \\
        q_Z
         & =
        A-B-C+D
        =
        2(A+D)-1.
    \end{aligned}
    \label{eq:pauli-channel-ptm}
\end{equation}
Here, \(q_X,q_Y,q_Z\) are also the eigenvalues of the single-qubit
Pauli channel on \(X,Y,Z\), respectively. It is obvious that $q_X,q_Y,q_Z\in [-1,1],
$
which leads to the identification challenge in QNT, we will elaborate this later in \S \ref{subsec:direct-nt}.

Consequently, the output state in
\eqref{eq:pauli-channel-output-state} has the Pauli decomposition
\begin{equation}
    \rho
    =
    \frac{1}{4}
    \left(
    I I
    +q_X X X
    -q_Y Y Y
    +q_Z Z Z
    \right),
    \label{eq:pauli-channel-output-expansion}
\end{equation}
and its restricted Pauli--Liouville vector is
\begin{equation}
    \plket{\rho}
    =
    \begin{bmatrix}
        1 & q_X & -q_Y & q_Z
    \end{bmatrix}^{T}.
    \label{eq:pauli-channel-output-vector}
\end{equation}


By Lemma~\ref{lem:bell-pauli-transformation}, the Pauli coefficients
of state $\rho$ uniquely determine the channel parameters
\(A,B,C,D\).


\subsection{Double-Selection Protocol}
\label{doub}\label{subsec:doub}

Entanglement distillation aims to  probabilistically
 transform multiple copies of a noisy
entangled pairs into  fewer pairs with higher fidelity to a target Bell
state, which we take to be $\ket{\Phi^+}$.
For example, the protocol considered here uses two noisy ancillary
pairs to probabilistically increase the fidelity of a third pair.
Recent work has shown that measurement statistics collected during
entanglement distillation can also be used for state
estimation~\cite{Casapao_2025}.

We consider the double-selection protocol introduced by Fujii and
Yamamoto~\cite{PhysRevA.80.042308}.
Figure~\ref{fig:double-selection-circuit} illustrates its dual role:
successful attempts yield purified entangled pairs, while measurement
records from all attempts are used to estimate the input
Bell-diagonal state under the conditions specified below.


\begin{figure*}[tb]
    \centering
    \resizebox{\textwidth}{!}{%
        \begin{tikzpicture}[x=1cm,y=1cm,line cap=round,line join=round]
% ------------------------------------------------------------------
% Local styles and helper commands (kept inside the figure scope)
% ------------------------------------------------------------------
\tikzset{
  dsPanel/.style={draw=black!55,line width=0.8pt},
  dsArrow/.style={-{Triangle[length=4.5mm,width=4.2mm]},draw=black!55,line width=3.2pt},
  dsWire/.style={draw=black,line width=0.75pt},
  dsEnt/.style={decorate,decoration={snake,amplitude=0.82mm,segment length=2.7mm},draw=black,line width=0.72pt},
  dsNoise/.style={decorate,decoration={snake,amplitude=1.55mm,segment length=5.2mm},draw=black!32,line width=0.65pt}
}

% Helper: glossy qubit node
\newcommand{\dsBall}[3]{%
  \shade[ball color=#3,draw=black!65,line width=0.55pt] (#1,#2) circle (2.35mm);
  \fill[white,opacity=0.34] (#1-0.045,#2+0.05) circle (0.55mm);
}

% Helper: smaller glossy qubit node
\newcommand{\dsMiniBall}[3]{%
  \shade[ball color=#3,draw=none] (#1,#2) circle (1.75mm);
  \fill[white,opacity=0.30] (#1-0.03,#2+0.032) circle (0.38mm);
}

% Helper: Bell pair between two horizontal stations
% A single sinusoidal-looking strand on a straight baseline.
\newcommand{\dsPair}[3]{%
  \draw[
    decorate,
    decoration={snake,amplitude=0.95mm,segment length=3.3mm},
    draw=black,
    line width=0.82pt
  ] (#1+0.15,#3) -- (#2-0.15,#3);
  \dsBall{#1}{#3}{blue!22!violet!45}
  \dsBall{#2}{#3}{blue!22!violet!45}
}

% Helper: a single sinusoidal-looking vertical entanglement strand
\newcommand{\dsVerticalEnt}[3]{%
  \draw[
    decorate,
    decoration={snake,amplitude=0.78mm,segment length=3.0mm},
    draw=black,
    line width=0.82pt
  ] (#1,#2) -- (#1,#3);
}

% Helper: CNOT symbols
\newcommand{\dsCtrl}[2]{\fill (#1,#2) circle (1.05mm);}
\newcommand{\dsTarg}[2]{%
  \draw[line width=0.8pt] (#1,#2) circle (2.25mm);
  \draw[line width=0.8pt] (#1-0.16,#2)--(#1+0.16,#2);
  \draw[line width=0.8pt] (#1,#2-0.16)--(#1,#2+0.16);
}

% Helper: meter box
\newcommand{\dsMeter}[3]{%
  \draw[line width=0.72pt] (#1-0.28,#2-0.22) rectangle (#1+0.28,#2+0.22);
  \draw[line width=0.65pt] (#1-0.16,#2-0.08) arc[start angle=180,end angle=0,radius=0.16];
  \draw[line width=0.65pt] (#1,#2-0.08)--(#1+0.085,#2+0.055);
  \node[anchor=west,font=\scriptsize] at (#1+0.36,#2) {$#3$};
}

% ------------------------------------------------------------------
% Left panel: noisy entangled pairs
% ------------------------------------------------------------------
\draw[dsPanel] (0.15,0.85) rectangle (4.25,6.95);
\dsPair{0.82}{3.58}{6.02}
\dsPair{0.82}{3.58}{5.18}
\dsPair{0.82}{3.58}{4.34}
\dsPair{0.82}{3.58}{3.50}
\fill (2.20,2.98) circle (0.045);
\fill (2.20,2.78) circle (0.045);
\fill (2.20,2.58) circle (0.045);
\node[font=\bfseries\normalsize] at (2.20,1.75) {Noisy entangled pairs};
\node[font=\normalsize] at (2.20,1.30) {Fidelity: $F_{\mathrm{in}}=0.80$};

% Arrow into the protocol
\draw[dsArrow] (4.27,4.02)--(4.92,4.02);

% ------------------------------------------------------------------
% Middle panel: double-selection protocol
% ------------------------------------------------------------------
\draw[draw=blue!75!black,line width=1.1pt,rounded corners=2.8mm]
      (4.95,1.18) rectangle (11.60,7.28);
\node[font=\bfseries\Large,text=blue!75!black] at (8.28,6.90)
      {Double-selection protocol};

% Wire y-coordinates
\def\ya{6.12}\def\yb{5.42}\def\yc{4.72}
\def\yA{3.42}\def\yB{2.72}\def\yC{2.02}
% staggered input x-coordinates
\def\xa{5.45}\def\xb{5.83}\def\xc{6.21}
\def\xend{10.62}

% Alice wires
\draw[dsWire] (\xa,\ya)--(\xend,\ya);
\draw[dsWire] (\xb,\yb)--(\xend,\yb);
\draw[dsWire] (\xc,\yc)--(\xend,\yc);
% Bob wires
\draw[dsWire] (\xa,\yA)--(\xend,\yA);
\draw[dsWire] (\xb,\yB)--(\xend,\yB);
\draw[dsWire] (\xc,\yC)--(\xend,\yC);

% Three shared Bell pairs: curved vertical wavy lines
\dsVerticalEnt{\xa}{\ya-0.13}{\yA+0.13}
\dsVerticalEnt{\xb}{\yb-0.13}{\yB+0.13}
\dsVerticalEnt{\xc}{\yc-0.13}{\yC+0.13}
\foreach \xx/\yt/\ybott in {\xa/\ya/\yA,\xb/\yb/\yB,\xc/\yc/\yC}{
  \dsMiniBall{\xx}{\yt}{blue!28!violet!42}
  \dsMiniBall{\xx}{\ybott}{blue!28!violet!42}
}

% Alice local double selection
\dsCtrl{8.02}{\ya}\dsTarg{8.02}{\yb}\draw[line width=0.8pt](8.02,\ya)--(8.02,\yb);
\dsTarg{9.18}{\yb}\dsCtrl{9.18}{\yc}\draw[line width=0.8pt](9.18,\yb)--(9.18,\yc);
\dsMeter{10.52}{\yb}{Z}
\dsMeter{10.52}{\yc}{X}

% Bob local double selection
\dsCtrl{8.02}{\yA}\dsTarg{8.02}{\yB}\draw[line width=0.8pt](8.02,\yA)--(8.02,\yB);
\dsTarg{9.18}{\yB}\dsCtrl{9.18}{\yC}\draw[line width=0.8pt](9.18,\yB)--(9.18,\yC);
\dsMeter{10.52}{\yB}{Z}
\dsMeter{10.52}{\yC}{X}

% Branching arrows
\draw[dsArrow] (11.60,5.62)--(12.38,6.15);
\draw[dsArrow] (11.60,2.88)--(12.38,2.54);

% ------------------------------------------------------------------
% Upper-right-middle panel: distilled copies
% ------------------------------------------------------------------
\draw[dsPanel] (12.40,4.20) rectangle (17.12,7.28);
\dsPair{13.42}{16.10}{6.79}
\dsPair{13.42}{16.10}{6.15}
\fill (14.76,5.80) circle (0.04);
\fill (14.76,5.62) circle (0.04);
\fill (14.76,5.44) circle (0.04);
\node[font=\bfseries\normalsize] at (14.76,5.08) {Distilled copies};
\node[font=\normalsize] at (14.76,4.57) {Fidelity: $F_{\mathrm{out}}\approx0.89$};

% ------------------------------------------------------------------
% Lower-right-middle panel: measurement statistics
% ------------------------------------------------------------------
\draw[dsPanel] (12.40,0.85) rectangle (17.12,3.46);
% Three histogram icons
\foreach \x/\shade in {13.20/45,14.76/72,16.32/92}{
  \fill[blue!\shade,rounded corners=1.1mm] (\x-0.50,2.20) rectangle (\x+0.50,3.04);
  \fill[white] (\x-0.34,2.36) rectangle (\x-0.23,2.82);
  \fill[white] (\x-0.12,2.36) rectangle (\x-0.01,2.54);
  \fill[white] (\x+0.10,2.36) rectangle (\x+0.21,2.69);
  \fill[white] (\x+0.32,2.36) rectangle (\x+0.41,2.88);
  \draw[white,line width=0.55pt] (\x-0.37,2.32)--(\x+0.43,2.32);
}
\node[font=\bfseries\normalsize] at (14.76,1.15) {Measurement statistics};
% \node[font=\normalsize,align=center] at (14.76,1.91)
%       {$\hat p_X$\\[-1mm]$\hat p_Z$\\[-1mm]$\hat p_{XZ}$};
\node[font=\normalsize,align=center] at (13.20,1.91)
      {$\hat p_X$};

      \node[font=\normalsize,align=center] at (14.76,1.91)
      {$\hat p_Z$};
      \node[font=\normalsize,align=center] at (16.32,1.91)
      {$\hat p_{XZ}$};
\draw[dsArrow] (17.12,2.62)--(17.91,2.62);

% ------------------------------------------------------------------
% Right panel: state estimation
% ------------------------------------------------------------------
\draw[dsPanel] (17.94,0.85) rectangle (24.58,6.98);
\node[font=\bfseries\Large] at (21.26,6.52) {State estimation};

% Axes
\draw[-{Triangle[length=2.8mm,width=2.6mm]},line width=1.0pt]
      (18.82,1.83)--(24.25,1.83);
\draw[-{Triangle[length=2.8mm,width=2.6mm]},line width=1.0pt]
      (18.82,1.83)--(18.82,6.08);

% Confidence band
\fill[blue!18,opacity=0.70]
  plot[smooth] coordinates {
    (18.88,5.72)(19.05,4.65)(19.22,3.85)(19.42,3.20)(19.62,2.78)
    (19.84,2.48)(20.05,2.30)(20.29,2.18)(20.53,2.10)(20.77,2.05)
    (21.04,2.01)(21.34,1.98)(21.66,1.96)(22.02,1.95)(22.40,1.94)
    (22.82,1.93)(23.24,1.92)(23.65,1.91)(24.02,1.90)}
  --
  plot[smooth] coordinates {
    (24.02,1.84)(23.65,1.845)(23.24,1.85)(22.82,1.855)(22.40,1.86)
    (22.02,1.87)(21.66,1.88)(21.34,1.89)(21.04,1.91)(20.77,1.94)
    (20.53,1.97)(20.29,2.02)(20.05,2.10)(19.84,2.22)(19.62,2.40)
    (19.42,2.72)(19.22,3.15)(19.05,3.75)(18.88,4.65)} -- cycle;

% Mean RMSE curve
\draw[blue!80!black,line width=1.15pt]
  plot[smooth] coordinates {
    (18.88,5.10)(19.05,4.20)(19.22,3.45)(19.42,2.90)(19.62,2.55)
    (19.84,2.30)(20.05,2.16)(20.29,2.08)(20.53,2.02)(20.77,1.98)
    (21.04,1.95)(21.34,1.93)(21.66,1.915)(22.02,1.905)(22.40,1.895)
    (22.82,1.888)(23.24,1.882)(23.65,1.876)(24.02,1.87)};
\node[font=\normalsize] at (21.55,1.43) {Repetitions $N$};
\node[font=\normalsize,rotate=90] at (18.34,4.00) {Estimation error (9MSE)};


% In-figure note, following the reference graphic
\end{tikzpicture}
%
    }
    \caption{
        Illustration of the double-selection protocol, with an example
        demonstrating its dual role in entanglement distillation and
        state estimation.
        Initially, Alice and Bob share independent copies of a Bell-diagonal
        state with fidelity $0.80$ with respect to $\ket{\Phi^+}$
        and equal weights on the other three Bell states.
        The subsystems $A_i$ and $B_i$ form the $i$th shared entangled pair.
        In each attempt, three pairs undergo bilateral CNOT operations,
        after which the second and third pairs are measured in the
        $Z$ and $X$ bases, respectively.
        Under ideal local operations and measurements, the first pair is
        retained when Alice's and Bob's outcomes coincide in both bases,
        yielding an output state with fidelity approximately $0.89$.
        Meanwhile, measurement records from all attempts, including
        unsuccessful ones, are used to estimate the input Bell-diagonal state.
    }
    \label{fig:double-selection-circuit}
\end{figure*}


Alice and Bob initially share three independent copies of the
Bell-diagonal state (called \emph{source state}) in \eqref{eq:pauli-channel-output-expansion} through a noisy quantum channel. Denote the copies by
\(\rho_1,\rho_2,\rho_3\), with
\(\rho_1=\rho_2=\rho_3=\rho\).

The protocol proceeds as follows.
\begin{enumerate}
    \item Alice and Bob each locally apply a  CNOT gate, using their qubit
          from $\rho_1$ as the control and their qubit from $\rho_2$ as
          the target.

    \item They each locally apply  a second local CNOT gate, using their qubit
          from $\rho_3$ as the control and their qubit from $\rho_2$ as
          the target.
    \item They measure \(\rho_2\) in the \(Z\) basis and \(\rho_3\) in
          the \(X\) basis. Let \(z_A,z_B\) and \(x_A,x_B\) denote the local
          outcomes of Alice and Bob, respectively.

    \item They retain \(\rho_1\) if and only if
          \(z_A=z_B\) and \(x_A=x_B\); otherwise, they discard it.
\end{enumerate}


We repeat the double-selection circuit and comparing Alice's and Bob's
measurement outcomes yields the marginal and joint coincidence
probabilities
\begin{equation}
    \begin{aligned}
        p_Z
         & \coloneqq \Pr(z_A=z_B)
        =\frac{1+q_Z^3}{2},       \\
        p_X
         & \coloneqq \Pr(x_A=x_B)
        =\frac{1+q_X^2}{2},       \\
        p_{XZ}
         & \coloneqq
        \Pr(z_A=z_B,\;x_A=x_B)=\frac{1+q_Z^3+q_X^2+q_Zq_Y^2}{4}.
    \end{aligned}
    \label{eq:double-selection-probabilities}
\end{equation}
Here, \(p_{XZ}\) is also the \emph{acceptance probability} of the protocol.

The following estimation result follows from the double-selection statistics derived by Casapao et al.~\citep[Sec.~III-A]{Casapao_2025}.

\begin{lemma}[Estimation from double-selection statistics]
    \label{lem:double-selection-statistics}
    Let \(\rho\) be the Bell-diagonal state in
    \eqref{eq:pauli-channel-output-expansion}. If
    \(F(\rho)>1/2\), then, under ideal local operations and measurements,
    the coincidence probabilities \(p_X,p_Z,p_{XZ}\) uniquely determine
    the nontrivial PTM eigenvalues \(q_X,q_Y,q_Z\).
    Specially, from rearranging~\eqref{eq:double-selection-probabilities}, we have
    \begin{equation}
        \begin{aligned}
            q_Z
             & =\sqrt[3]{2p_Z-1}, \\
            q_X
             & =\sqrt{2p_X-1},    \\
            q_Y
             & =
            \sqrt{
                \frac{
                    4p_{XZ}-2p_X-2p_Z+1
                }{
                    \sqrt[3]{2p_Z-1}
                }
            }.
        \end{aligned}
        \label{eq:double-selection-inversion}
    \end{equation}
    Together with
    \eqref{eq:pauli-channel-output-vector} and
    Lemma~\ref{lem:bell-pauli-transformation}, these quantities uniquely
    determine the input Bell-diagonal state.
\end{lemma}

% \begin{proof}
% The condition \(F(\rho)>1/2\) implies
% \(q_X,q_Y,q_Z>0\). Therefore, the positive solution of
% \eqref{eq:double-selection-probabilities} is unique and is given by
% Hence, \(p_X,p_Z,p_{XZ}\) uniquely determine
% \(q_X,q_Y,q_Z\). 
% \end{proof}



    \begin{observation}
        \label{lem3}
        The estimation of parameter $q_Z$ in the double-selection protocol is unique by itself and does not rely on the fidelity condition \(F(\rho)> \frac 1 2\).
    \end{observation}
    Later on, we will leverage this observation and another channel dressing technique to remove the requirement of the \(F(\rho)> \frac 1 2\) condition in our protocols.



